%! Author = sriadityavedantam
%! Date = 4/8/26

\section{Linear Systems of Divisors}

Let $X$ be a nonsingular variety.
All Weil divisors are Cartier on nonsingular varieties.
If we have $D \subseteq X$ a divisor, there is an associated vector space $\Ell(D)$ defined by
\[
    \Ell(D) = \{0\} \cup \{f \in \c(X)\setminus\{0\} : \div(f) + D \geq 0\}.
\]
If
\[
    \div(f) + D = \sum a_i C_i,
\]
then $a_i \geq 0$ for all $i$.
Thus $\div(f) + D$ is effective.
If $D = \sum n_i C_i$, then $f \in \Ell(D)$ if and only if $\nu_{C_i}(f) \geq -n_i$ for each $i$, and $\nu_C(f) \geq 0$ for any $C \neq C_i$.
Since $D$ is Cartier, there exists $\o(D)$, a line bundle, and
\[
    \Ell(D) = H^0(X,\o(D)),
\]
where $H^0$ denotes the space of sections of $\o(D)$.
If $X = \pp^1$ and $D = nX_{\infty}$, where the point at infinity is $[1:0]$, then we ask: what is $\Ell(D)$?
Let $f \in \c(\pp^1)\setminus\{0\}$.
Then $f = p/q$, where $p,q$ are homogeneous polynomials of the same degree, say $d$.
Let $\{p_1,\ldots,p_d\}$ denote the zeros of $p(x_0,x_1)$, and similarly for $q$.
Then
\[
    \div(f) = \sum p_i - \sum q_i.
\]
We want $\div(f) + D \geq 0$, where $D = nX_{\infty}$.
If there exists some zero $q_i \neq X_{\infty}$, then this is not true.
Otherwise,
\begin{align*}
    \div(f) + D &= \sum p_i - dX_{\infty} + nX_{\infty} \\
    &= \sum p_i + (n-d)X_{\infty}.
\end{align*}
This is effective if and only if $n-d \geq 0$.
Then $q = cX_1^d$, and we have
\[
    f = \frac{\sum_{i=0}^d a_i X_0^i X_1^{d-i}}{cX_1^d}.
\]
If $a_i' = a_i/c$, then
\[
    f = \sum_{i=0}^d a_i' \left(\frac{x_0}{x_1}\right)^i.
\]
So $f \in \c[x_0/x_1]_{\leq n}$, meaning the highest degree is at most $n$.
Hence $\Ell(D)$ is a $\c$-vector space of dimension $n+1$.
If $n \leq -1$, then $\Ell(D) = \{0\}$.

\subsection{Hypersurfaces}

\begin{theorem*}{
    If $X$ is a smooth projective variety, then $\Ell(D)$ is finite-dimensional.
}
\end{theorem*}

We write $\dim_\c(\Ell(D)) = \ell(D)$.
Why do we need projectivity?
Let $D$ be the zero divisor.
Then
\[
    \Ell(D) = \{0\} \cup \{f \in \c(X)\setminus\{0\} : \div(f) \geq 0\}.
\]
So $f$ has no poles, hence $f$ is regular.
If $X = \pp^1$, then $\Ell(D) = \c[\pp^1] = \c$.
If $X = \aa^1$, then $\Ell(D) = \c[\aa^1] = \c[t]$.
Then $\ell(D) = \infty$.

\begin{definition}[Linear System]
    Let $X$ be projective.
    Then the linear system of $D \in \Div(X)$ is the projectivization of $\Ell(D)$, denoted $\pp(\Ell(D))$.
    This is isomorphic to
    \[
        \frac{\Ell(D)\setminus\{0\}}{\c^\upstar} = \pp^{\ell(D)-1}.
    \]
\end{definition}

\begin{proposition}{
    Suppose $X$ is projective.
    If $D \sim D^\prime$, then $\Ell(D) \cong \Ell(D^\prime)$.
}
    If $D \sim D^\prime$, then $D - D^\prime = \div(g)$ for some $g \in \c(X)$.
    Consider
    \[
        \varphi : \Ell(D) \to \Ell(D^\prime)
    \]
    defined by
    \[
        \varphi(f) = fg.
    \]
    If $f \in \Ell(D)$, then $\div(f) + D \geq 0$.
    Now
    \begin{align*}
        \div(fg) + D^\prime &= \div(f) + \div(g) + D^\prime \\
        &= \div(f) + (D - D^\prime) + D^\prime \\
        &= \div(f) + D \geq 0.
    \end{align*}
    So $fg \in \Ell(D^\prime)$.
    Conversely, the inverse map
    \[
        \varphi^{-1} : \Ell(D^\prime) \to \Ell(D)
    \]
    is defined by
    \[
        h \mapsto \frac{h}{g}.
    \]
    Thus $\Ell(D) \cong \Ell(D^\prime)$.
\end{proposition}

In particular, $\pp(\Ell(D)) \cong \pp(\Ell(D^\prime))$.
One can check that $\pp(\Ell(D))$ is naturally identified with the set of effective divisors linearly equivalent to $D$.
If $f \in \Ell(D)\setminus\{0\}$, then
\[
    \div(f) + D \geq 0
\]
is an effective divisor linearly equivalent to $D$.
Moreover, if $c \in \c^\upstar$, then $\div(cf) = \div(f)$, so the associated effective divisor is unchanged.
Let $X = \pp^n$ and let $D = \{f_d = 0\}$ be a hypersurface of degree $d$.
Let $D^\prime$ be defined similarly.
Then
\[
    D - D^\prime = \div\left(\frac{f_d}{f_d^\prime}\right),
\]
since $\deg(f_d) = \deg(f_d^\prime)$.
Hence
\[
    \Ell(D) \cong \{\text{degree } d \text{ homogeneous polynomials in } x_0,\ldots,x_n\}.
\]
Also,
\[
    \frac{f_d}{f_d^\prime} \in \c(\pp^n).
\]
We already know how many degree $d$ homogeneous polynomials there are, by stars and bars.
Therefore
\[
    \pp(\Ell(D)) = \pp^{\binom{n+d}{d}-1},
\]
which is the projective space of all hypersurfaces of degree $d$ in $\pp^n$.
Let $D$ be a cubic curve in $\pp^2$.
Then
\[
    \pp(\Ell(D)) \cong \pp^9.
\]

We will later show that for a genus $g$ curve, the divisor class group is not countable.
With Riemann--Roch, we know $Cl(\pp^n) = \z$, while $Cl\!\left(\sum p\right)$ is not countable.
This is not subtle.